\subsubsection{Mixed historical--synthetic pretraining}
The terminal curriculum experiment pooled 100 unique vine-synthetic episodes and 61 historical-prefix episodes into 1,000 pretraining presentations before the unchanged 61-episode historical fine-tuning stage.  On the 22 complete holding periods, the mixed ensemble attained the highest CAGR and annualized CRRA certainty equivalent among the four matched training protocols, while using materially less turnover and implementation drag than historical-only training.  Historical-only training nevertheless retained the lowest volatility, drawdown, and CVaR and the highest Sharpe, Sortino, Calmar, and Omega ratios.

The registered mixed-minus-control CE effects were positive in the primary panel, but none survived Holm adjustment.  In particular, the mixed-minus-historical ensemble effect was not representative of matched seeds: fewer than half of the seedwise effects were positive and the median effect was negative.  Nevertheless, the ensemble result was not generated by one indispensable seed: eight of ten leave-one-seed-out mixed ensembles retained a positive CE difference relative to historical-only training, and all ten retained positive differences relative to both synthetic controls.  The strongest directional evidence was against synthetic-only training, showing that synthetic experience did not substitute for real historical information.  The results therefore support synthetic simulation as an augmentation and implementation regularizer rather than a standalone training distribution.  Because the experiment was designed after inspection of the consumed holdout, it is reported as post-holdout explanatory evidence and the mixed curriculum is retained only as a candidate for independently frozen external validation.
